% ===================================================================
%  نقشہ نمبر 16 — وڈی دکان۔ — منڈی۔ — کھڈوݨے۔
%  Traduction 2026 d'apres texto/io, controlee sur texto/fr, texto/en
%  et texto/hi. Consigne : voir texto/pnb/10-naqsha-01.tex.
%
%  Le tableau d'astronomie (bloc 15-1) porte des renvois a LETTRES,
%  (a) a (f), graves sur la planche elle-meme : ils se rangent sous
%  le numero 85, comme le dit gravuri/literi.json. L'ido y saute la
%  parenthese ouvrante du (d) ; la traduction la donne, et le
%  controle passe ce bloc sans rien dire.
%
%  LES TROIS MASQUES GARDENT LEUR NOM, comme le chien Medor et l'ane
%  Aliboron : ہارلیکن, پولیشینیل, پیرو. Ce sont des personnages, non
%  des objets.
% ===================================================================

%%K t16-apar-1 apar
\VUsaut{4.0mm}
\VUtitre{15.0pt}{60}{نقشہ نمبر 16}
\VUsaut{2.0mm}
\VUpk{13.2pt}{40}{وڈی دکان۔ --- منڈی۔}
\VUpk{13.2pt}{40}{کھڈوݨے۔}
\VUsaut{3.4mm}

%%K t16-01-1 p
1. --- نواں ورھا آ رہیا اے۔ وڈیاں دکاناں نے \VUgras{کھڈوݨیاں} \textsuperscript{(1)}
تے نویں ورھے دیاں سوغاتاں دی اپݨی سجاوٹ تیار کر لئی اے۔

%%K t16-01-2 p
میں اپݨے دادا نوں آکھیا پئی مینوں ایہہ سوہݨی سجاوٹ ویکھݨ منڈی لے چلو، تے اوہ من
گئے۔

%%K t16-02-1 p
2. --- پہلاں اسیں ہتھیاراں دیاں کجھ سوہݨیاں سجاوٹاں دے سامݨے رُکے، جیہناں وچ اک
\VUgras{رائفل}، اک \VUgras{جھالر والی ٹوپی}، اک \VUgras{کِرچ}، اک
\VUgras{تلوار دی پیٹی} تے اک جوڑی \VUgras{مونڈھے دے پھُندنے} سن، یا فیر اک
\VUgras{خود}، اک \VUgras{چھاتی دا کوچ} \textsuperscript{(3)} تے اک
\VUgras{پستول} \textsuperscript{(4)}۔ ایہہ سب مینوں \VUgras{جادوئی لالٹیناں}
\textsuperscript{(2)} نالوں کِتے ودھ بھایا۔

%%K t16-tit-1 sub
\VUsaut{3.0mm}
\VUpk{10.2pt}{40}{ویکھݨ جوگ چیزاں۔}
\VUsaut{2.6mm}

%%K t16-03-1 p
3. --- اوسے شعبے وچ اک \VUgras{پہرے دار} \textsuperscript{(5)} اپݨی
\VUgras{پہرے دی کوٹھڑی} وچ سی؛ اوہ گتے دے پورے چڑیا گھر اُتے پہرہ دیندا لگدا سی۔
کِنّے جانور سن اوتھے! اپݨے اڈے اُتے اک \VUgras{طوطا} \textsuperscript{(6)}، نک اُتے
سِنگ لئی اک \VUgras{گینڈا} \textsuperscript{(7)}، اک \VUgras{بارہ سنگھا}
\textsuperscript{(8)}، اک \VUgras{شتر مرغ} \textsuperscript{(9)}، اک \VUgras{بندر}
\textsuperscript{(10)} جیہڑا اخروٹ توڑ رہیا سی، اک \VUgras{لومڑی}
\textsuperscript{(11)} جیہڑی اک کُکڑی چُکی لے جاندی سی، اک \VUgras{مگرمچھ}
\textsuperscript{(12)} اپݨے وڈے جبڑے کھولی، اک \VUgras{اُٹھ} \textsuperscript{(13)}،
اک سجیلا \VUgras{گرے ہاؤنڈ} \textsuperscript{(14)}، اک \VUgras{تیندوا}
\textsuperscript{(15)}، تے فیر دو جانور جیہناں نوں میں جاݨدا نہ سی، جیہڑے، مینوں
دسیا گیا، اک \VUgras{نیولا} \textsuperscript{(16)} تے اک \VUgras{بھیڑیا}
\textsuperscript{(17)} نیں۔ اوتھے اک \VUgras{چیتا} \textsuperscript{(18)} تے اک
\VUgras{بگھیاڑ} \textsuperscript{(21)} وی سی؛ کول ای پیارے نِکے \VUgras{چوہے}
\textsuperscript{(19)} تے ربڑ دیاں نِکیاں \VUgras{چُوہیاں} \textsuperscript{(20)}
سن۔ چھیکر اک \VUgras{دریائی گھوڑا} \textsuperscript{(22)} تے اک کُب والا
\VUgras{اُٹھ} \textsuperscript{(23)} سنگرہ پورا کردے سن۔ میں اوتھے گھنٹیاں کھڑا
رہندا، جے کول ای \VUgras{ٹین دے سپاہیاں} \textsuperscript{(29)} نال بھریا اک شاندار
پڑاء میری اکھ نہ کھِچ لیندا۔

%%K t16-tit-2 sub
\VUsaut{4.0mm}
\VUpk{10.2pt}{40}{سپاہی تے جنگ۔}
\VUsaut{2.8mm}

%%K t16-04-1 p
4. --- میرے دادا، جیہڑے \VUgras{اپاہج سابق سپاہی} \textsuperscript{(24)} نیں تے
جیہناں نوں اک \VUgras{زخم} مگروں فوج چھڈݨی پئی، جیہنے اوہناں نوں
\VUgras{فوجی تمغہ} دلوایا، اوہناں نوں مینوں اوہناں \VUgras{مہماں} دیاں گلاں سُݨاؤݨا
بہت چنگا لگدا اے جیہناں وچ اوہ شامل ہوئے۔ اوہ مینوں اوس نِکی فوج دیاں وکھو وکھ
چالاں سمجھا کے بہت خوش ہوئے۔ منڈی ویکھݨ آئے اصلی سپاہیاں دا اک ٹولہ — اک
\VUgras{انجینیر سپاہی} \textsuperscript{(25)}، اک \VUgras{پہاڑی ہولا سپاہی}
\textsuperscript{(26)} اپݨی \VUgras{ٹوپی} وچ، پیدل فوج دا اک \VUgras{مکھ حوالدار}
\textsuperscript{(27)} تے توپ خانے دا اک \VUgras{حوالدار} \textsuperscript{(28)}،
جیدی آستین اُتے سونے دی \VUgras{دھاری} سی — اوہ وی وضاحت سُݨن ساڈے کول رُک گئے۔

%%K t16-05-1 p
5. --- میرے دادا، اِنّے مشہور سُݨن والے لبھ کے بہت مان نال بھرے، اوہناں جھٹ جنگ دی
اک پوری وِیؤنت گھڑ سُٹی۔

%%K t16-05-2 p
اوہ قلعہ بند شہر، اپݨیاں \VUgras{فصیلاں} تے \VUgras{کھائیاں} سمیت،
\VUgras{ویری فوج} نال گھِریا سی، جیہڑی لمی \VUgras{گولہ باری} مگروں اوہدے اُتے
\VUgras{دھاوا} بولݨ دی تیاری وچ سی۔ پر \VUgras{گھِرے ہوئے} لوک، بھُکھ توں مجبور،
\VUgras{گھیرن والیاں} دے \VUgras{پڑاء} اُتے اک \VUgras{باہر واہ} کر بیٹھے سن۔
\VUgras{تنبواں} دے چوہاں پاسیں اک تِکھی \VUgras{لڑائی} وچ \VUgras{ہلے} نوں پچھے دھکّ
کے \VUgras{جِتّݨ والے} گھیرن والیاں نے اپݨے \VUgras{گھوڑ سوار دستے}
\VUgras{ہارے ہوئے} حملہ آوراں دا \VUgras{پچھا} کرن گھلے۔ اوہ \VUgras{پچھے ہٹے}،
\VUgras{رݨ بھوں} نوں \VUgras{مُردیاں} تے \VUgras{زخمیاں} نال وِچھا کے۔
\VUgras{اسٹریچر والے} زخمیاں نوں \VUgras{چلدے دواخانے} تیکر لے گئے، تاں جو
\VUgras{جراح} اوہناں دی سنبھال کرن۔

%%K t16-05-3 p
کجھ \VUgras{بٹالینواں} دے سورمے مقابلے دے باوجود، جیہناں نے \VUgras{چوکور ویوہ}
بݨایا سی، \VUgras{ہار} پوری سی؛ بچی فوج \VUgras{نس گئی}، بہت سارے \VUgras{بندی}
پچھے چھڈ کے۔ ویری شہر لے کے اپݨی \VUgras{جِت} اُتے تاج رکھݨ نوں سی۔ شاید ایہو مہم دا
انت ہوندا، تے اک \VUgras{جنگ بندی} مگروں اک \VUgras{امن دے سمجھوتے} اُتے دستخط ہو
سکدے۔

%%K t16-tit-3 sub
\VUsaut{4.4mm}
\VUpk{10.2pt}{40}{نویں ورھے دیاں سوغاتاں۔}
\VUsaut{2.8mm}

%%K t16-06-1 p
6. --- اوہ جوان سپاہی، جیہڑے سابق سپاہی دی رنگیلی کہاݨی سُݨدے ہوئے ہس رہے سن،
اوہناں نے اوہناں کولوں کجھ ہور سوال پُچھے۔ رہیا میں، تے میں کاؤنٹر دا چکر کٹ کے اک
سوہݨی \VUgras{کمان} \textsuperscript{(32)} تے اک \VUgras{دھنُش} \textsuperscript{(33)}
اوہدے \VUgras{تیر} \textsuperscript{(31)} سمیت ویکھݨ گیا۔ اوتھے مینوں جانوراں دا اک
ہور سنگرہ مِلیا: دو \VUgras{گاوݨ والیاں چڑیاں} \textsuperscript{(34)}، اک چابی والی
\VUgras{بلبل} تے \VUgras{پھُدکی} جیہڑیاں گلا پھاڑ کے گا رہیاں سن، تے انوکھے جیواں دی
اک قطار — اک \VUgras{چمگادڑ} \textsuperscript{(35)}، اک \VUgras{اجگر}
\textsuperscript{(36)}، اک \VUgras{کچھو} \textsuperscript{(37)}، اک \VUgras{سیل}
\textsuperscript{(38)}، اک \VUgras{وھیل} \textsuperscript{(39)} تے سونے دی جھِلی نال
بݨی اک \VUgras{شارک} \textsuperscript{(40)}۔

%%K t16-07-1 p
7. --- میں اوہناں نوں ویکھݨ ودھ دیر نہ رُکیا، کیوں جے میری اکھ اوس شے اُتے پے چُکی
سی جیدا میں سُفنا ویکھدا سی: اک شاندار \VUgras{کٹھ پُتلی رنگ منچ}
\textsuperscript{(41)}، شاندار \VUgras{منظر پٹ} تے اک سوہݨے لال \VUgras{پردے} سمیت۔
\VUgras{منچ} اُتے دو اداکار کِسے بہت دلچسپ \VUgras{ناٹک} وچ اپݨی \VUgras{بھومکا}
نبھاندے لگدے سن، کیوں جے ڈبیاں دے گتے دے نِکے \VUgras{ویکھݨ والے}، یا جیہڑے
\VUgras{اگلیاں کرسیاں} تے \VUgras{پچھلے حصے} وچ \VUgras{ساز چالک} دے پچھے بیٹھے سن،
اوہ گرم جوشی نال تالی وجا رہے سن۔

%%K t16-07-2 p
افسوس، اوہ رنگ منچ بہت مہنگا سی۔

%%K t16-08-1 p
8. --- ایہدے توں بغیر، کھڈوݨیاں وچوں چوݨا اوکھا سی، کیوں جے کھڑکیاں ہر طرحاں دے
کھڈوݨیاں نال بھریاں سن: \VUgras{ہارلیکن} \textsuperscript{(42)}،
\VUgras{پولیشینیل} \textsuperscript{(43)}، \VUgras{پیرو} \textsuperscript{(44)}، پر
پھڑ پھڑاندے \VUgras{اُلو} \textsuperscript{(45)}، سیٹی وجاندے \VUgras{کالکنٹھ}
\textsuperscript{(46)}، بھونکدے \VUgras{پوڈل}، اک \VUgras{گرامو فون}
\textsuperscript{(47)} تے \VUgras{بُجھارتاں}۔ ایہناں وچوں اک کھڈوݨا مینوں بہت بھایا:
اوہدے وچ اک \VUgras{سمندری لڑائی} \textsuperscript{(48)} وکھائی گئی سی، جیدے وچ
\VUgras{ناریل دے رُکھاں} نال بھرے کنڈھے دے سامݨے \VUgras{سمندری ڈاکو} اک
\VUgras{جہاز} اُتے ہلہ کردے نیں۔

%%K t16-noto-1 noto
\VUnotes{143.2mm}{%
\textit{(*) نقشہ نمبر 6 ویکھو۔}%
}

%%K t16-09-1 p
9. --- میں ایہہ ساریاں حیران کر دیݨ والیاں چیزاں بڑے سکون نال ویکھ سکیا، کیوں جے
دکان وچ لوک گھٹ ای سن — مشکل نال مُٹھ بھر تماشبین، اک \VUgras{گھوڑ سوار سپاہی}
\textsuperscript{(49)}، تے اک \VUgras{وصولی کرن والا} \textsuperscript{(50)}، جیہڑا
مکھ \VUgras{گلّے} \textsuperscript{(51)} اُتے اک \VUgras{ہنڈی} پیش کر رہیا سی۔

%%K t16-10-1 p
10. --- میں نے اپݨے دادا کولوں پُچھیا پئی \VUgras{خزانچی} دے پچھے طاقاں اُتے رکھیاں
اوہ کتاباں کیس کم دیاں نیں؛ اوہناں دسیا پئی اوہ \VUgras{حساب دیاں وہیاں} نیں۔

%%K t16-tit-4 sub
\VUsaut{3.6mm}
\VUpk{10.2pt}{40}{دکان وچ جھاتی۔}
\VUsaut{2.8mm}

%%K t16-11-1 p
11. --- کھڈوݨیاں دا شعبہ چھڈ کے اسیں زین سازی ول گئے، تاں جو \VUgras{زیناں}
\textsuperscript{(53)}، \VUgras{رکاباں} \textsuperscript{(54)}، \VUgras{لگاماں}
\textsuperscript{(55)}، \VUgras{دہاݨیاں} \textsuperscript{(56)} تے
\VUgras{سواری دے چابکاں} اُتے اک نظر مار لئیے۔ جدوں تیکر \VUgras{شعبے دے منتظم}
\textsuperscript{(57)} میرے دادا نوں کجھ جاݨکاری دیندے رہے، میں \VUgras{چینی مِٹی}
تے \VUgras{شیشے دے برتناں} \textsuperscript{(58)} دی سجاوٹ ویکھݨ گیا، جِتھے سوہݨے
\VUgras{سلاد دے کٹورے} تے نازک \VUgras{چا دان} \textsuperscript{(59)} سن۔

%%K t16-12-1 p
12. --- پہلی منزل اُتے چڑھن لئی اسیں \VUgras{کمر بنداں} \textsuperscript{(60)} دی
کھڑکی دے پچھوں لنگھے، جیہڑی جُرابا دے شعبے کول اے، جِتھے بہت سوہݨیاں \VUgras{نِکراں}
\textsuperscript{(63)} وِچھیاں سن۔ کول ای اک \VUgras{دکان دی مددگار}
\textsuperscript{(61)} اک \VUgras{گاہک} \textsuperscript{(62)} نوں سوہݨے ریشمی
\VUgras{کپڑے} \textsuperscript{(64)} چوݨ وچ مدد کر رہی سی۔

%%K t16-12-2 p
پوڑیاں چڑھدے ہوئے میں ویکھیا پئی دکان جاپانی \VUgras{لالٹیناں} \textsuperscript{(65)}،
\VUgras{چھتریاں} \textsuperscript{(66)} تے وڈیاں \VUgras{کھجوراں} \textsuperscript{(70)}
نال سجی اے۔

%%K t16-13-1 p
13. --- پہلی منزل اُتے ویکھݨ نوں ودھ کجھ نہ سی: نِرا ساز سامان (*)،
\VUgras{تجوریاں} \textsuperscript{(67)}، \VUgras{درازاں والیاں الماریاں}
\textsuperscript{(68)} تے \VUgras{بال گڈیاں} \textsuperscript{(69)}۔ پر دکان دا سارا
منظر بہت \VUgras{پرچاوا دیݨ والا} سی۔

%%K t16-14-1 p
14. --- اپݨے تھلے مینوں ساز وکھائی دِتے: \VUgras{وِیݨاواں} \textsuperscript{(71)}،
\VUgras{گٹار} \textsuperscript{(72)}، \VUgras{مینڈولن} \textsuperscript{(73)}،
\VUgras{کارنیٹ} \textsuperscript{(74)}، \VUgras{کلیرینیٹ} \textsuperscript{(75)}،
اپݨے \VUgras{گز} \textsuperscript{(76)} سمیت وائلن تے اک \VUgras{وڈا ڈھول}
\textsuperscript{(77)} تیکر۔ کجھ ہور اگے، لکھݨ دے سامان دے کاؤنٹر اُتے، اک
\VUgras{طالب علم} \textsuperscript{(78)} \VUgras{دکان دے مددگار} \textsuperscript{(79)}
نال اک کاغذ دے تھیلے دا مُل بھاء کر رہیا سی۔ اوہناں کول مینوں اک سوہݨی
\VUgras{قاعدے دی کتاب} \textsuperscript{(80)} وکھائی دِتی۔

%%K t16-14-2 p
دکان دے وچکار اک اُچا \VUgras{کرسمس دا رُکھ} \textsuperscript{(81)} کھڑا سی، جیدے
اُتے چوہاں پاسیں موم بتیاں، گہݨے تے ہر طرحاں دے کھڈوݨے ٹنگے سن:
\VUgras{لکڑ دے گھوڑے} \textsuperscript{(82)}، \VUgras{گُڈیاں} \textsuperscript{(83)}،
وغیرہ۔

%%K t16-14-3 p
\VUgras{عطر دی دکان} \textsuperscript{(84)}، اپݨے \VUgras{دند دے منجناں} تے اپݨے وکھو
وکھ \VUgras{عطراں} سمیت، اک بہت سوہݨی وا چھڈدی سی جیہڑی ساڈے تیکر اُٹھ آندی سی۔

%%K t16-tit-5 sub
\VUsaut{3.4mm}
\VUpk{10.2pt}{40}{اسیں کجھ خریدنے آں۔}
\VUsaut{2.8mm}

%%K t16-15-1 p
15. --- کیوں جے میرے دادا نوں ویلا لما لگݨ لگ پیا سی، اسیں وڈی پوڑی توں مُڑ تھلے
اُترے۔ اوہ اک \VUgras{تارا ودیا دی سُوچی} \textsuperscript{(85)} دے سامݨے پل بھر رُکے،
جیدے اُتے، اوہناں دسیا، \VUgras{دُم دار تارے} \textsuperscript{(a)}،
\VUgras{چن تے سورج دے گرہݨ} \textsuperscript{(b)}، \VUgras{چوہاں سیدھاں}
\textsuperscript{(c)} \VUgras{(اُتر، دکھݨ، چڑھدا تے لہندا)} والا اک سیدھ چکر، دوہاں
\VUgras{اَدھ گولیاں} \textsuperscript{(d)} دا نقشہ، \VUgras{گرہ} \textsuperscript{(e)}
تے \VUgras{شمسی نظام} \textsuperscript{(f)} وکھائے گئے سن۔ مینوں اوہدے وچوں کجھ سمجھ
نہ آیا، تے مینوں کِتے ودھ دلچسپی اوس \VUgras{سامان اپڑاؤݨ والے} \textsuperscript{(86)}
دی گوٹے والی وردی وچ ہوئی جیہڑا لنگھ رہیا سی، تے اپݨے کول رکھیاں اوہناں سوہݨیاں
\VUgras{پکھیاں} \textsuperscript{(87)} وچ، \VUgras{بٹویاں} \textsuperscript{(88)} تے
\VUgras{جیبی بٹویاں} \textsuperscript{(89)} کول۔

%%K t16-16-1 p
16. --- میں نے کاؤنٹر اُتے کجھ \VUgras{پر} \textsuperscript{(90)} تے
\VUgras{پیٹی دے بکسوے} \textsuperscript{(91)} وی تکے، تے اوہ سوہݨے
\VUgras{بݨاوٹی پھُل} \textsuperscript{(94)} جیہڑے ٹوکریاں وچ سلیقے نال سجے سن۔
\VUgras{بنفشے}، \VUgras{گُل شب بو}، \VUgras{سنبل} \textsuperscript{(95)}،
\VUgras{جیرینیم} \textsuperscript{(96)}، \VUgras{کنول} \textsuperscript{(97)} تے
\VUgras{نسترن} \textsuperscript{(98)} دی نقل بہت چنگی سی۔

%%K t16-tit-6 sub
\VUsaut{4.4mm}
\VUpk{10.2pt}{40}{دادا دی سوغات۔}
\VUsaut{2.8mm}

%%K t16-17-1 p
17. --- میں نے اپݨے دادا نوں آکھیا پئی مینوں اوہناں \VUgras{دند دے برشاں}
\textsuperscript{(92)} وچوں اک خرید دیو جیہڑے \VUgras{والاں دیاں پِناں}
\textsuperscript{(93)} کول اک ڈبے وچ سن۔ اوہ من گئے، تے میں آپ گلّے اُتے اپݨے خریدے
دا مُل تارن گیا، جیدے کول اک وڈا \VUgras{پیکٹ} \textsuperscript{(100)} بنھیا جا رہیا
سی، اک \VUgras{گاہک} \textsuperscript{(99)} لئی۔

%%K t16-18-1 p
18. --- اچانک کلاتمک \VUgras{کانسی دے بُتاں} \textsuperscript{(101)} کول کھڑے لوکاں
وچ ہلچل مچ گئی۔ اک \VUgras{چور} \textsuperscript{(102)} نوں ہُݨے ہُݨے اک
\VUgras{دکان دے نگران} \textsuperscript{(103)} نے کِسے دی جیب کٹدے رنگے ہتھیں پھڑ لیا
سی۔ اوہنوں \VUgras{گرفتار} کرن نوں سپاہی سدیا گیا، تے اسیں نکل ای رہے سی پئی مجرم
نوں تھاݨے لے جایا جا رہیا سی۔
\VUsaut{6.0mm}
\VUfilet{22mm}
